\chapter{Etapa 1 del proyecto financiado}
\label{ch:stage1}
\label{ch:wbs_stage1}

La metodología doctoral requiere una transición controlada desde evidencia
computacional hacia mediciones físicas y, después, hacia un sitio real. La Etapa 1
del proyecto financiado proporciona esa infraestructura. Este capítulo conserva
en la monografía sólo las dependencias científicas, los productos verificables y
los criterios que habilitan la etapa siguiente; cronogramas semanales, compras,
movilidad y logística se mantienen en la documentación operativa del repositorio.

\section{Finalidad y relación con la tesis}

Durante los primeros diez meses, la etapa integra modelado, diseño experimental,
plataforma multimodal y preparación del sitio comunitario. Su éxito no se mide por
volumen de datos ni número de participantes, sino por reducción demostrable de
incertidumbre: un modelo reproducible, un observable caracterizado, una plataforma
estable, un protocolo trazable y condiciones justificadas para avanzar.

La evidencia estática disponible se presenta en el \cref{ch:avances}. Su consecuencia
para el proyecto es concreta: la siguiente incertidumbre central concierne a la
sensibilidad física. E0 mide el piso instrumental y E1 compara el Jacobiano simulado
con el producido por un movimiento conocido. Las auditorías computacionales y la
comparación de representaciones pueden continuar en paralelo, pero no sustituyen
esas intervenciones. Del mismo modo, un banco estable no acredita todavía sensado
humano ni habilita por sí solo un despliegue comunitario.

\section{Líneas de trabajo y productos científicos}

Las seis líneas de trabajo se conectan por dependencia, no sólo por calendario.
LT1 formula las predicciones; LT2 establece qué información debe conservarse; LT3
materializa la medición; LT4 somete el sistema a perturbaciones controladas; LT5
caracteriza el dominio de despliegue; y LT6 integra la evidencia para decidir si
procede la Etapa 2.

\begin{table}[htbp]
\centering\small
\caption{Líneas de trabajo, función metodológica y evidencia de cierre.}
\label{tab:wbs_stage1_summary}
\begin{tabularx}{\textwidth}{P{1.0cm}P{3.0cm}L L}
\toprule
Línea & Función & Producto principal & Condición de cierre\\
\midrule
LT1 & Modelo físico y gemelo inicial & Modelo directo y discrepancia cuantificada &
Reproducción y evaluación fuera de calibración.\\
LT2 & Datos, protocolo y referencias & Protocolo, metadatos y control de calidad &
Trazabilidad y reglas fijadas antes de la campaña.\\
LT3 & Plataforma multimodal & Banco E0--E1 y sincronización & Piso, deriva,
referencia y geometría caracterizados.\\
LT4 & Estudio controlado & Evidencia acumulativa E0--E3 & Cada régimen satisface
su criterio antes de aumentar complejidad.\\
LT5 & Sitio y vinculación & Ficha técnica y caso de uso & Contraparte,
restricciones y RF del sitio documentados.\\
LT6 & Integración & Paquete de transición & Riesgos y supuestos pendientes
explícitos.\\
\bottomrule
\end{tabularx}
\end{table}

LT1 y LT2 preservan la distinción entre modelo, observable y tarea. La simulación
debe producir $H$ bajo convenciones reproducibles; el protocolo debe especificar
cómo se obtiene $Y$ y qué referencias quedan disponibles en inferencia. Los productos
incluyen el modelo directo, la comparación con datos públicos, el esquema multimodal,
los procedimientos de inclusión y exclusión y los métodos de referencia sin selección
oráculo.

LT3 implementa la arquitectura funcional del \cref{ch:platform}. El banco mínimo
combina radio, desplazamiento mecánico medido, geometría registrada y marcas
temporales. Radar, RGB-D, fisiología e inerciales son modalidades auxiliares según
la pregunta; su presencia en el inventario no implica que se utilicen en toda captura.
La estabilidad se demuestra mediante E0 y no mediante especificaciones nominales.

LT4 sigue la jerarquía E0--E3. La habitación vacía establece falsas perturbaciones;
el reflector impone posición y movimiento; el fantoma introduce forma de onda
respiratoria; y una persona añade deformación, postura y variabilidad biológica con
la aprobación ética correspondiente. La comparación Wi-Fi/LoRa utiliza, cuando sea
posible, la misma perturbación y referencia para separar diferencias del operador
de diferencias de estado.

LT5 no constituye validación clínica. Su función científica es describir un cambio
de dominio antes de medirlo: geometría, materiales, energía, infraestructura, ruido
RF, cobertura y restricciones de montaje. El taller de divulgación y el codiseño
ambiental pertenecen a vinculación; no se cuentan como observaciones de sujetos.
LT6 reúne estos productos y conserva tanto resultados favorables como condiciones
que impidan avanzar.

\section{Dependencia experimental y criterios de revisión}

La secuencia de evidencia es
\[
\begin{aligned}
\text{modelo y protocolo}&\rightarrow \text{E0: instrumento}
\rightarrow \text{E1: movimiento}\\
&\rightarrow \text{E2: fantoma}
\rightarrow \text{E3: persona}
\rightarrow \text{sitio}.
\end{aligned}
\]
Un cumplimiento cronológico no reemplaza una condición de paso. Las revisiones
se resumen en la \cref{tab:stage1_milestones}; sus umbrales cuantitativos se fijan
después del piloto técnico y antes de la evaluación confirmatoria, tal como establece
el \cref{ch:methodology}.

\begin{table}[htbp]
\centering\small
\caption{Revisiones de la Etapa 1 y decisión que sustentan.}
\label{tab:stage1_milestones}
\begin{tabularx}{\textwidth}{P{1.4cm}L L}
\toprule
Revisión & Evidencia requerida & Decisión\\
\midrule
R1 & Modelo estático reproducible; discrepancia y acceso a fase declarados. &
Fijar el modelo físico de referencia.\\
R2 & Protocolo, calidad, gobernanza y ética formalizados. &
Autorizar adquisición humana cuando corresponda.\\
R3 & E0--E1 con referencia, sincronización e incertidumbre caracterizadas. &
Continuar hacia el fantoma.\\
R4 & E0--E2 con repetibilidad, sensibilidad y recuperación de forma de onda. &
Iniciar E3 y preparar transferencia.\\
R5 & Sitio, contraparte, restricciones y caso de uso definidos. &
Preparar instalación posterior.\\
$R_{\rm E2}$ & Cumplimiento conjunto de todos los bloques. &
Iniciar la campaña formal de Etapa 2.\\
\bottomrule
\end{tabularx}
\end{table}

La condición conjunta es
\begin{equation}
R_{\rm E2}=
R_{\rm RF}\land
R_{\rm sinc}\land
R_{\rm protocolo}\land
R_{\rm etica}\land
R_{\rm sitio}\land
R_{\rm datos}.
\label{eq:stage2_gate}
\end{equation}
La conjunción impide compensar una referencia inestable con mayor volumen de datos
o sustituir aprobación ética por madurez técnica. Si una condición falla, el
producto de la etapa es la delimitación documentada de ese bloqueo y la revisión
del diseño correspondiente.

\section{Recursos como restricciones del diseño}

La instrumentación y el presupuesto condicionan qué observable puede adquirirse,
pero no constituyen aportes científicos. La prioridad técnica es resolver referencia
RF, movimiento medido, geometría y sincronización antes de ampliar modalidades. Se
reutilizan las plataformas disponibles y cada compra se justifica por una incertidumbre
experimental: acceso a I/Q, coherencia, precisión mecánica, registro o cadencia.

El detalle de inventario, bolsas autorizadas, cotizaciones por validar, movilidad,
kits, cronograma semanal y actividades del taller se conserva en
\path{docs/etapa1_detalle_operativo.md}. Separarlo del argumento principal evita
que una decisión administrativa se lea como evidencia, sin borrar la trazabilidad
necesaria para ejecutar el proyecto.

\section{Transición}

La Etapa 1 conecta la teoría con una campaña auditable. El gemelo predice qué cambio
debería observarse; el banco establece qué cambio puede medirse; E1 y E2 contrastan
si ambas respuestas concuerdan; y el sitio define el dominio al que se intentará
transferir. El capítulo final sintetiza qué se conoce ya, qué permanece como hipótesis
y qué evidencia tiene mayor valor para continuar la investigación.
